%!TEX root = ../thesis.tex

\chapter{Future Work}\label{chap:future-work}

There are various possible follow-ups to this work that are concerned with expanding
the single-seed base case, improving the EfficientZero learner, modifying the fuzzing
loop, and including additional benchmark suites and ablations.

\section{Full seed corpus}

Since the goal is fuzzing stateful applications, we want to extend the basic
single-seed corpus approach we have used throughout this thesis to a full seed
corpus example keeping track of multiple fixed seeds at once as described
in~\cref{sec:corpus-level-actions}.

This includes evaluating retain and select actions and incorporating new
information into the observation model to enable the agent to learn a good
policy.

\section{Improving the EfficientZero learner}

We have presented the MuZero lineage in~\cref{sec:ez-lineage} up to EfficientZero
and chose it as the base of our implementation.
EfficientZero itself was released in 2021 and is therefore somewhat dated.
There are various successors that introduce additional improvements and sample efficiency
gains.
Incorporating these improvements into the base EfficientZero architecure we
implemented here is an interesting line of future work.

Multiple prior works exist that attack the problem of action space size explosion (\cref{eq:action-space-size})
we are likely to encounter when implementing the full corpus-input model.
Sampled MuZero~\cite{hubert_learning_2021} samples a subset of actions and plans over them
instead of the whole action space while correcting the policy-improvement operator for the sampling distribution.
Gumbel MuZero~\cite{danihelka_policy_2022} guarantees policy improvement by sampling
actions without replacement using the Gumbel-Top-k trick~\cite{kool_stochastic_2019}.
Multiagent Gumbel MuZero~\cite{hao_multiagent_2024} tackles the action space explosion
problem by modeling it as a multi-agent problem.
EfficientZero V2~\cite{wang_efficientzero_2024} combines the results of Sampled
MuZero and Gumbel MuZero.

Other works improve the world model and the throughput.
SpeedyZero~\cite{mei_speedyzero_2023} improves the raw execution performance of EfficientZero
by introducing a distributed and highly parallel architecture.
UniZero~\cite{pu_unizero_2025} uses a transformer as the world model to improve long-term memory
and reduce problems caused by partial observability.
ReZero~\cite{xuan_rezero_2024} modifies the reanalyze stage to improve execution performance.
LightZero~\cite{niu_lightzero_2023} contains reference implementations of ReZero and UniZero
and evaluates them via a unified benchmark.

\section{Modifying the fuzzing loop}

In our implementation, the learner is strictly \gls{gpu}-bound.
This leads to a waste of resources since the \gls{cpu} is underutilized.
Ideally, we want to do useful \gls{cpu} work while waiting for the \gls{gpu} to complete
its operation.
One possible approach is a hybrid fuzzing scheme: by running a classical \gls{cgf} fuzzer
alongside our RL-based fuzzer, we can make full use of all system resources.
The different fuzzers can then exchange information, improving overall performance.

\section{Additional benchmark suites and ablations}

We have tested our implementation against a few small examples, proving its feasibility.
Next, we want to include more and larger targets to measure exactly how parameters
like the \gls{mcts} simulation count impacts the fuzzing of targets of higher complexity.
Additionally, we expect further ablation studies and a detailed parameter tuning to improve fuzzing
performance.
